
\documentclass{article}
\usepackage{colortbl}
\usepackage{makecell}
\usepackage{multirow}
\usepackage{supertabular}

\begin{document}

\newcounter{utterance}

\centering \large Interaction Transcript for game `cladder', experiment `full\_v1.5\_default', episode 1489 with refineonly{-}Qwen3.5{-}9B.
\vspace{24pt}

{ \footnotesize  \setcounter{utterance}{1}
\setlength{\tabcolsep}{0pt}
\begin{supertabular}{c@{$\;$}|p{.15\linewidth}@{}p{.15\linewidth}p{.15\linewidth}p{.15\linewidth}p{.15\linewidth}p{.15\linewidth}}
    \# & \multicolumn{2}{c}{Player} && \multicolumn{2}{c}{Game Master} \\
    \hline

    \theutterance \stepcounter{utterance}  
    & & & \multicolumn{4}{p{0.6\linewidth}}{
        \cellcolor[rgb]{0.9,0.9,0.9}{
            \makecell[{{p{\linewidth}}}]{
                \texttt{\tiny{[P1$\langle$GM]}}
                \texttt{You are an expert in causal inference. The following question is not a typical commonsense query, but rather a meticulously designed question created by a professor specializing in causal inference, intended to assess the students' mastery of the course content.} \\
\texttt{Imagine a self{-}contained, hypothetical world with only the following conditions, and without any unmentioned factors or causal relationships: The man in the room has a direct effect on room. The candle has a direct effect on room. For people not blowing out candles, the probability of dark room is 88\%. For people who blow out candles, the probability of dark room is 56\%. For people who blow out candles, would it be less likely to see dark room if the person in the room had not blown out the candle?} \\
\texttt{Start your answer with 'Yes' or 'No', followed by additional reasoning or evidence to support your explanation.} \\
            }
        }
    }
     \\ \\

    \theutterance \stepcounter{utterance}  
    & \multicolumn{4}{p{0.6\linewidth}}{
        \cellcolor[rgb]{0.561,0.737,0.561}{
            \makecell[{{p{\linewidth}}}]{
                \texttt{\tiny{[P1$\rangle$GM]}}
                \texttt{No.} \\
\\ 
\texttt{**Reasoning:**} \\
\\ 
\texttt{To answer this, we must analyze the causal structure and the specific counterfactual query using the provided probabilities.} \\
\\ 
\texttt{1.  **Identify the Variables and Probabilities:**} \\
\texttt{    *   Let \$B\$ be the action of blowing out the candle (\$B=1\$) or not (\$B=0\$).} \\
\texttt{    *   Let \$D\$ be the state of the room being dark (\$D=1\$) or not (\$D=0\$).} \\
\texttt{    *   The prompt provides the following conditional probabilities:} \\
\texttt{        *   \$P(D=1 | B=0) = 0.88\$ (Probability of a dark room if the candle is *not* blown out).} \\
\texttt{        *   \$P(D=1 | B=1) = 0.56\$ (Probability of a dark room if the candle *is* blown out).} \\
\\ 
\texttt{2.  **Analyze the Causal Mechanism:**} \\
\texttt{    *   The prompt states the candle has a "direct effect on room." In this context, blowing out a candle (action \$B=1\$) reduces the probability of the room being dark compared to not blowing it out (\$0.56 < 0.88\$). This implies that in this specific hypothetical world, the act of blowing out the candle causes the room to be less likely to be dark (perhaps the "blowing out" action triggers a light switch, or the data simply defines the causal direction where \$B=1\$ leads to lower darkness).} \\
\texttt{    *   Crucially, the data establishes that \$B=1\$ is the condition associated with a *lower* probability of darkness (\$56\textbackslash{}\%\$) compared to \$B=0\$ (\$88\textbackslash{}\%\$).} \\
\\ 
\texttt{3.  **Evaluate the Counterfactual Question:**} \\
\texttt{    *   **Scenario:** "For people who blow out candles" (Actual state: \$B=1\$).} \\
\texttt{    *   **Question:** "Would it be less likely to see dark room if the person in the room had **not** blown out the candle?"} \\
\texttt{    *   **Translation:** We are comparing the probability of darkness in the actual world (\$B=1\$, where \$P(D=1) = 0.56\$) against the counterfactual world where the person did not blow out the candle (\$B=0\$, where \$P(D=1) = 0.88\$).} \\
\texttt{    *   **Comparison:** The question asks if the probability would be *less* likely in the counterfactual scenario (\$B=0\$).} \\
\texttt{    *   **Calculation:** Is \$P(D=1 | B=0)\$ less than \$P(D=1 | B=1)\$?} \\
\texttt{    *   \$0.88 < 0.56\$ is **False**.} \\
\texttt{    *   In fact} \\
            }
        }
    }
    & & \\ \\

    \theutterance \stepcounter{utterance}  
    & & & \multicolumn{2}{p{0.3\linewidth}}{
        \cellcolor[rgb]{0.9,0.9,0.9}{
            \makecell[{{p{\linewidth}}}]{
                \texttt{\tiny{[GM$|$GM]}}
                \texttt{No} \\
            }
        }
    }
    & & \\ \\

    \theutterance \stepcounter{utterance}  
    & & & \multicolumn{2}{p{0.3\linewidth}}{
        \cellcolor[rgb]{0.9,0.9,0.9}{
            \makecell[{{p{\linewidth}}}]{
                \texttt{\tiny{[GM$|$GM]}}
                \texttt{no} \\
            }
        }
    }
    & & \\ \\

    \theutterance \stepcounter{utterance}  
    & & & \multicolumn{2}{p{0.3\linewidth}}{
        \cellcolor[rgb]{0.9,0.9,0.9}{
            \makecell[{{p{\linewidth}}}]{
                \texttt{\tiny{[GM$|$GM]}}
                \texttt{game\_result = WIN} \\
            }
        }
    }
    & & \\ \\

\end{supertabular}
}

\end{document}
